
\mdef   \label{exa:keyn5la21} Consider the $J_{1,5}$-module $\Deltaaa{{5}}{({3},(21))}$. By choosing a basis $\{\nC_{(2,1)}, (2 \ 3) \nC_{(2,1)}\}$ for $\SS_{(2,1)}$, we can give the following diagram basis for this module where the rectangular box ``receiving''  %``absorbing" 
%\ppm{[-why absorbing? across?]}\bajm{[because the lines go in and don't come out?]}\ppm{[receiving?]}\bajm{[if you are desparate to change it then do so, but I don't think any of the options you suggested are obviously preferable, but perhaps I'm missing something :)]}
three propagating lines denotes the idempotent $\nC_{(2,1)}$. 
%The lex-ordered basis given by the diagram order \ppm{[i guess this actually not quite right yet!...]} is:  
\begin{align}
&\hspace{-1cm} 
\onecuppppp{6}{1}{2}{3}{4}{5}  \; 
\onecuppppp{6}{1}{3}{2}{4}{5}  \;
\onecuppppp{6}{1}{4}{2}{3}{5}  \;
\onecuppppp{6}{2}{3}{1}{4}{5} \;
\onecuppppp{6}{2}{4}{1}{3}{5}  \;
\onecuppppp{6}{3}{4}{1}{2}{5}  \;
\onecuppppp{6}{4}{5}{1}{2}{3}  
\nn\\
&\hspace{1cm} 
\onecupppppx{6}{1}{2}{3}{4}{5}  \;
\onecupppppx{6}{1}{3}{2}{4}{5}    \;
\onecupppppx{6}{1}{4}{2}{3}{5}  \;
\onecupppppx{6}{2}{3}{1}{4}{5}  \;  
\onecupppppx{6}{2}{4}{1}{3}{5}  \;
\onecupppppx{6}{3}{4}{1}{2}{5}  \;
\onecupppppx{6}{4}{5}{1}{2}{3}  \label{eq:21basis}
\end{align}
The value of the form on two diagram basis elements is determined by applying * to one of the diagrams (flipping it upside down), vertically juxtaposing the diagrams, and computing the associated scalar. Schematically then, we compute 
\begin{align*}&\reflectbox{\rotatebox[]{180}{\onecuppppp{6}{1}{2}{3}{4}{5}}}=\a,
&&\reflectbox{\rotatebox[]{180}{\onecuppppp{6}{1}{2}{3}{4}{5}}}=1.\\[-0.822em]
&\ \onecuppppp{6}{1}{2}{3}{4}{5}
&&\ \onecuppppp{6}{1}{3}{2}{4}{5}\end{align*}

\noindent 
and so on. In particular 
\ignore{{
\bajm{[I am now convinced we actually need to switch to the EFE convention for the symmetriser - this is because the explicit matrices we are writing will be wrong otherwise... for example, when $\l=(2,1)$, computing $\langle (12),e\rangle$ gives 1 in the EFE formalism since $(1 2) E=E$, yet it gives $-1$ in the FE formalism since $(12)F=-F$... For example, using $C_\l$ as written here we find  ]}
\begin{align*}&\reflectbox{\rotatebox[]{180}{\onecuppppp{6}{1}{2}{3}{4}{5}}}=C_{\l}^* (12) C_{\l}=-C_\l^* C_\l \Rightarrow \left\langle \onecuppppp{6}{1}{2}{3}{4}{5} , \onecuppppp{6}{1}{4}{2}{3}{5} \right\rangle =-1
\\[-1.7em]
&\ \onecuppppp{6}{1}{4}{2}{3}{5}\end{align*}
\bajm{[However, this should correspond to the (1,4) entry of the Gram matrix below for $\Deltaaa{{5}}{({3},(21))}$, which is 1... ]}\bajm{[I am now proceeding with the updated convention since this gives the correct matrix!]}\ppm{[agreed.  (I haven't got my head around the diff yet, but your way is not wrong.)]}
\medskip 
}}
\begin{align*}& \reflectbox{\rotatebox[]{180}{\onecuppppp{6}{1}{4}{2}{3}{5}}}=-\tfrac{1}{2},
&&\reflectbox{\rotatebox[]{180}{\onecuppppp{6}{1}{4}{2}{3}{5}}}=-\tfrac{1}{2},
&&\reflectbox{\rotatebox[]{180}{\onecuppppp{6}{1}{4}{2}{3}{5}}}=0,
\\[-0.822em]
& \ \onecupppppx{6}{1}{3}{2}{4}{5}
&&  \ \onecupppppx{6}{2}{4}{1}{3}{5}
&& \ \onecupppppx{6}{2}{3}{1}{4}{5}\\
&\reflectbox{\rotatebox[]{180}{\onecuppppp{6}{1}{4}{2}{3}{5}}}=-\tfrac{1}{2},
&&\reflectbox{\rotatebox[]{180}{\onecuppppp{6}{1}{4}{2}{3}{5}}}=-\tfrac{1}{2}.
\\[-0.822em]
& \ \onecupppppx{6}{3}{4}{1}{2}{5}
&& \ \onecupppppx{6}{4}{5}{1}{2}{3}
\end{align*}
\noindent 
%Alternatively we can use an orthonormal basis for the Specht part, such as in (\ref{eq:21ortho}). 
We give the Gram matrix for $\Deltaaa{{5}}{({3},(21))}$ below, where we use the orthonormal basis as per \ref{eq:21ortho} for $\SS_{(2,1)}$ (in the same order as \eqref{eq:21basis} but with $(2 \ 3) \nC_{(2,1)}$ replaced by $\frac{2}{\sqrt{3}}((2 \ 3)+\frac{1}{2}e )\nC_{(2,1)}$): 
%\bajm{[maybe give it without the lines??]}

\[
\left(
\begin{array}{ccccccc ccccccc}
 \alpha  & 1 & 1 & 1 & 1 & 0 & 0 & 0 & 0 & 0 & 0 & 0 & 0 & 0 \\
 1 & \alpha  & 1 & 1 & 0 & 1 & 0 & 0 & 0 & 0 & 0 & 0 & 0 & 0 \\
 1 & 1 & \alpha  & 0 & 1 & 1 & -\frac{1}{2} & 0 & 0 & 0 & 0 & 0 & 0 & -\frac{\sqrt{3}}{2} \\
 1 & 1 & 0 & \alpha  & 1 & 1 & 0 & 0 & 0 & 0 & 0 & 0 & 0 & 0 \\
 1 & 0 & 1 & 1 & \alpha  & 1 & -\frac{1}{2} & 0 & 0 & 0 & 0 & 0 & 0 & \frac{\sqrt{3}}{2} \\
 0 & 1 & 1 & 1 & 1 & \alpha  & 1 & 0 & 0 & 0 & 0 & 0 & 0 & 0 \\
 0 & 0 & -\frac{1}{2} & 0 & -\frac{1}{2} & 1 & \alpha  & 0 & 0 & \frac{\sqrt{3}}{2} & 0 & \frac{\sqrt{3}}{2} & 0 & 0 
 \\ %\hline 
 0 & 0 & 0 & 0 & 0 & 0 & 0 & \alpha  & 1 & -1 & 1 & -1 & 0 & 0 \\
 0 & 0 & 0 & 0 & 0 & 0 & 0 & 1 & \alpha  & 1 & 1 & 0 & -1 & 0 \\
 0 & 0 & 0 & 0 & 0 & 0 & \frac{\sqrt{3}}{2} & -1 & 1 & \alpha  & 0 & 1 & -1 & -\frac{1}{2} \\
 0 & 0 & 0 & 0 & 0 & 0 & 0 & 1 & 1 & 0 & \alpha  & 1 & 1 & 0 \\
 0 & 0 & 0 & 0 & 0 & 0 & \frac{\sqrt{3}}{2} & -1 & 0 & 1 & 1 & \alpha  & 1 & \frac{1}{2} \\
 0 & 0 & 0 & 0 & 0 & 0 & 0 & 0 & -1 & -1 & 1 & 1 & \alpha  & 1 \\
 0 & 0 & -\frac{\sqrt{3}}{2} & 0 & \frac{\sqrt{3}}{2} & 0 & 0 & 0 & 0 & -\frac{1}{2} & 0 & \frac{1}{2} & 1 & \alpha 
\end{array}
\right)
\]
By direct computation, we find:
\begin{align*}\Delttt{5}{3,(21)}&= (\alpha -2)^3 \alpha  (\alpha +2) (\alpha +4)(\alpha ^4-7 \alpha ^2+3)^2,\\
\Delttt{6}{4,(21)}&=(\alpha -2)^3 \alpha  (\alpha +2) (\alpha +4)(\alpha -1)^2 \alpha^2(\alpha +1)^2 \left(\alpha ^2-7\right)^2.
\end{align*}


\ignore{\bajm{[-------------------------I would propose to end \S \ref{ss:sdmods} here-----------------------]}

Gram matrix for $\Deltaaa{7}{(5,(21))}$: 
\[
\left(
\begin{array}{ccccccc|cc|ccccccc|cc}
 \alpha  & 1 & 1 & 1 & 1 & 0 & 0 & 0 & 0 & 0 & 0 & 0 & 0 & 0 & 0 & 0 & 0 & 0 \\
 1 & \alpha  & 1 & 1 & 0 & 1 & 0 & 0 & 0 & 0 & 0 & 0 & 0 & 0 & 0 & 0 & 0 & 0 \\
 1 & 1 & \alpha  & 0 & 1 & 1 & -\frac{1}{2} & 0 & 0 & 0 & 0 & 0 & 0 & 0 & 0 & -\frac{\sqrt{3}}{2} & 0 & 0 \\
 1 & 1 & 0 & \alpha  & 1 & 1 & 0 & 0 & 0 & 0 & 0 & 0 & 0 & 0 & 0 & 0 & 0 & 0 \\
 1 & 0 & 1 & 1 & \alpha  & 1 & -\frac{1}{2} & 0 & 0 & 0 & 0 & 0 & 0 & 0 & 0 & \frac{\sqrt{3}}{2} & 0 & 0 \\
 0 & 1 & 1 & 1 & 1 & \alpha  & 1 & 0 & 0 & 0 & 0 & 0 & 0 & 0 & 0 & 0 & 0 & 0 \\
 0 & 0 & -\frac{1}{2} & 0 & -\frac{1}{2} & 1 & \alpha  & 1 & 0 & 0 & 0 & \frac{\sqrt{3}}{2} & 0 & \frac{\sqrt{3}}{2} & 0 & 0 & 0 & 0 \\ \hline 
 0 & 0 & 0 & 0 & 0 & 0 & 1 & \alpha  & 1 & 0 & 0 & 0 & 0 & 0 & 0 & 0 & 0 & 0 \\
 0 & 0 & 0 & 0 & 0 & 0 & 0 & 1 & \alpha  & 0 & 0 & 0 & 0 & 0 & 0 & 0 & 0 & 0 \\ \hline 
 0 & 0 & 0 & 0 & 0 & 0 & 0 & 0 & 0 & \alpha  & 1 & -1 & 1 & -1 & 0 & 0 & 0 & 0 \\
 0 & 0 & 0 & 0 & 0 & 0 & 0 & 0 & 0 & 1 & \alpha  & 1 & 1 & 0 & -1 & 0 & 0 & 0 \\
 0 & 0 & 0 & 0 & 0 & 0 & \frac{\sqrt{3}}{2} & 0 & 0 & -1 & 1 & \alpha  & 0 & 1 & -1 & -\frac{1}{2} & 0 & 0 \\
 0 & 0 & 0 & 0 & 0 & 0 & 0 & 0 & 0 & 1 & 1 & 0 & \alpha  & 1 & 1 & 0 & 0 & 0 \\
 0 & 0 & 0 & 0 & 0 & 0 & \frac{\sqrt{3}}{2} & 0 & 0 & -1 & 0 & 1 & 1 & \alpha  & 1 & \frac{1}{2} & 0 & 0 \\
 0 & 0 & 0 & 0 & 0 & 0 & 0 & 0 & 0 & 0 & -1 & -1 & 1 & 1 & \alpha  & 1 & 0 & 0 \\ %\hline 
 0 & 0 & -\frac{\sqrt{3}}{2} & 0 & \frac{\sqrt{3}}{2} & 0 & 0 & 0 & 0 & 0 & 0 & -\frac{1}{2} & 0 & \frac{1}{2} & 1 & \alpha  & 1 & 0 \\ \hline 
 0 & 0 & 0 & 0 & 0 & 0 & 0 & 0 & 0 & 0 & 0 & 0 & 0 & 0 & 0 & 1 & \alpha  & 1 \\
 0 & 0 & 0 & 0 & 0 & 0 & 0 & 0 & 0 & 0 & 0 & 0 & 0 & 0 & 0 & 0 & 1 & \alpha  \\
\end{array}
\right)\]

Gram matrix for $\Deltaaa{8}{(6,(21))}$:
\[ \hspace{-1cm} 
\left(
\begin{array}{ccccccc|ccc|ccccccc|ccc}
 \alpha  & 1 & 1 & 1 & 1 & 0 & 0 & 0 & 0 & 0 & 0 & 0 & 0 & 0 & 0 & 0 & 0 & 0 & 0 & 0 \\
 1 & \alpha  & 1 & 1 & 0 & 1 & 0 & 0 & 0 & 0 & 0 & 0 & 0 & 0 & 0 & 0 & 0 & 0 & 0 & 0 \\
 1 & 1 & \alpha  & 0 & 1 & 1 & -\frac{1}{2} & 0 & 0 & 0 & 0 & 0 & 0 & 0 & 0 & 0 & -\frac{\sqrt{3}}{2} & 0 & 0 & 0 \\
 1 & 1 & 0 & \alpha  & 1 & 1 & 0 & 0 & 0 & 0 & 0 & 0 & 0 & 0 & 0 & 0 & 0 & 0 & 0 & 0 \\
 1 & 0 & 1 & 1 & \alpha  & 1 & -\frac{1}{2} & 0 & 0 & 0 & 0 & 0 & 0 & 0 & 0 & 0 & \frac{\sqrt{3}}{2} & 0 & 0 & 0 \\
 0 & 1 & 1 & 1 & 1 & \alpha  & 1 & 0 & 0 & 0 & 0 & 0 & 0 & 0 & 0 & 0 & 0 & 0 & 0 & 0 \\
 0 & 0 & -\frac{1}{2} & 0 & -\frac{1}{2} & 1 & \alpha  & 1 & 0 & 0 & 0 & 0 & \frac{\sqrt{3}}{2} & 0 & \frac{\sqrt{3}}{2} & 0 & 0 & 0 & 0 & 0 
 \\  \hline  
 0 & 0 & 0 & 0 & 0 & 0 & 1 & \alpha  & 1 & 0 & 0 & 0 & 0 & 0 & 0 & 0 & 0 & 0 & 0 & 0 
 \\
 0 & 0 & 0 & 0 & 0 & 0 & 0 & 1 & \alpha  & 1 & 0 & 0 & 0 & 0 & 0 & 0 & 0 & 0 & 0 & 0 
 \\
 0 & 0 & 0 & 0 & 0 & 0 & 0 & 0 & 1 & \alpha  & 0 & 0 & 0 & 0 & 0 & 0 & 0 & 0 & 0 & 0 
 \\  \hline 
 0 & 0 & 0 & 0 & 0 & 0 & 0 & 0 & 0 & 0 & \alpha  & 1 & -1 & 1 & -1 & 0 & 0 & 0 & 0 & 0 \\
 0 & 0 & 0 & 0 & 0 & 0 & 0 & 0 & 0 & 0 & 1 & \alpha  & 1 & 1 & 0 & -1 & 0 & 0 & 0 & 0 \\
 0 & 0 & 0 & 0 & 0 & 0 & \frac{\sqrt{3}}{2} & 0 & 0 & 0 & -1 & 1 & \alpha  & 0 & 1 & -1 & -\frac{1}{2} & 0 & 0 & 0 \\
 0 & 0 & 0 & 0 & 0 & 0 & 0 & 0 & 0 & 0 & 1 & 1 & 0 & \alpha  & 1 & 1 & 0 & 0 & 0 & 0 \\
 0 & 0 & 0 & 0 & 0 & 0 & \frac{\sqrt{3}}{2} & 0 & 0 & 0 & -1 & 0 & 1 & 1 & \alpha  & 1 & \frac{1}{2} & 0 & 0 & 0 \\
 0 & 0 & 0 & 0 & 0 & 0 & 0 & 0 & 0 & 0 & 0 & -1 & -1 & 1 & 1 & \alpha  & 1 & 0 & 0 & 0 \\
 0 & 0 & -\frac{\sqrt{3}}{2} & 0 & \frac{\sqrt{3}}{2} & 0 & 0 & 0 & 0 & 0 & 0 & 0 & -\frac{1}{2} & 0 & \frac{1}{2} & 1 & \alpha  & 1 & 0 & 0 
 \\ \hline 
 0 & 0 & 0 & 0 & 0 & 0 & 0 & 0 & 0 & 0 & 0 & 0 & 0 & 0 & 0 & 0 & 1 & \alpha  & 1 & 0 \\
 0 & 0 & 0 & 0 & 0 & 0 & 0 & 0 & 0 & 0 & 0 & 0 & 0 & 0 & 0 & 0 & 0 & 1 & \alpha  & 1 \\
 0 & 0 & 0 & 0 & 0 & 0 & 0 & 0 & 0 & 0 & 0 & 0 & 0 & 0 & 0 & 0 & 0 & 0 & 1 & \alpha  \\
\end{array}
\right)
\]

Gram matrix for $\Deltaaa{9}{(7,(21))}$:
\[
\hspace{-.2in}
\left(
\begin{array}{ccccccc|cccc|ccccccc|cccc}
 \alpha  & 1 & 1 & 1 & 1 & 0 & 0 & 0 & 0 & 0 & 0 & 0 & 0 & 0 & 0 & 0 & 0 & 0 & 0 & 0 & 0 & 0 \\
 1 & \alpha  & 1 & 1 & 0 & 1 & 0 & 0 & 0 & 0 & 0 & 0 & 0 & 0 & 0 & 0 & 0 & 0 & 0 & 0 & 0 & 0 \\
 1 & 1 & \alpha  & 0 & 1 & 1 & -\frac{1}{2} & 0 & 0 & 0 & 0 & 0 & 0 & 0 & 0 & 0 & 0 & -\frac{\sqrt{3}}{2} & 0 & 0 & 0 & 0 \\
 1 & 1 & 0 & \alpha  & 1 & 1 & 0 & 0 & 0 & 0 & 0 & 0 & 0 & 0 & 0 & 0 & 0 & 0 & 0 & 0 & 0 & 0 \\
 1 & 0 & 1 & 1 & \alpha  & 1 & -\frac{1}{2} & 0 & 0 & 0 & 0 & 0 & 0 & 0 & 0 & 0 & 0 & \frac{\sqrt{3}}{2} & 0 & 0 & 0 & 0 \\
 0 & 1 & 1 & 1 & 1 & \alpha  & 1 & 0 & 0 & 0 & 0 & 0 & 0 & 0 & 0 & 0 & 0 & 0 & 0 & 0 & 0 & 0 \\
 0 & 0 & -\frac{1}{2} & 0 & -\frac{1}{2} & 1 & \alpha  & 1 & 0 & 0 & 0 & 0 & 0 & \frac{\sqrt{3}}{2} & 0 & \frac{\sqrt{3}}{2} & 0 & 0 & 0 & 0 & 0 & 0 
 \\  \hline 
 0 & 0 & 0 & 0 & 0 & 0 & 1 & \alpha  & 1 & 0 & 0 & 0 & 0 & 0 & 0 & 0 & 0 & 0 & 0 & 0 & 0 & 0 \\
 0 & 0 & 0 & 0 & 0 & 0 & 0 & 1 & \alpha  & 1 & 0 & 0 & 0 & 0 & 0 & 0 & 0 & 0 & 0 & 0 & 0 & 0 \\
 0 & 0 & 0 & 0 & 0 & 0 & 0 & 0 & 1 & \alpha  & 1 & 0 & 0 & 0 & 0 & 0 & 0 & 0 & 0 & 0 & 0 & 0 \\
 0 & 0 & 0 & 0 & 0 & 0 & 0 & 0 & 0 & 1 & \alpha  & 0 & 0 & 0 & 0 & 0 & 0 & 0 & 0 & 0 & 0 & 0 
 \\ \hline 
 0 & 0 & 0 & 0 & 0 & 0 & 0 & 0 & 0 & 0 & 0 & \alpha  & 1 & -1 & 1 & -1 & 0 & 0 & 0 & 0 & 0 & 0 \\
 0 & 0 & 0 & 0 & 0 & 0 & 0 & 0 & 0 & 0 & 0 & 1 & \alpha  & 1 & 1 & 0 & -1 & 0 & 0 & 0 & 0 & 0 \\
 0 & 0 & 0 & 0 & 0 & 0 & \frac{\sqrt{3}}{2} & 0 & 0 & 0 & 0 & -1 & 1 & \alpha  & 0 & 1 & -1 & -\frac{1}{2} & 0 & 0 & 0 & 0 \\
 0 & 0 & 0 & 0 & 0 & 0 & 0 & 0 & 0 & 0 & 0 & 1 & 1 & 0 & \alpha  & 1 & 1 & 0 & 0 & 0 & 0 & 0 \\
 0 & 0 & 0 & 0 & 0 & 0 & \frac{\sqrt{3}}{2} & 0 & 0 & 0 & 0 & -1 & 0 & 1 & 1 & \alpha  & 1 & \frac{1}{2} & 0 & 0 & 0 & 0 \\
 0 & 0 & 0 & 0 & 0 & 0 & 0 & 0 & 0 & 0 & 0 & 0 & -1 & -1 & 1 & 1 & \alpha  & 1 & 0 & 0 & 0 & 0 \\
 0 & 0 & -\frac{\sqrt{3}}{2} & 0 & \frac{\sqrt{3}}{2} & 0 & 0 & 0 & 0 & 0 & 0 & 0 & 0 & -\frac{1}{2} & 0 & \frac{1}{2} & 1 & \alpha  & 1 & 0 & 0 & 0 
 \\ \hline  
 0 & 0 & 0 & 0 & 0 & 0 & 0 & 0 & 0 & 0 & 0 & 0 & 0 & 0 & 0 & 0 & 0 & 1 & \alpha  & 1 & 0 & 0 \\
 0 & 0 & 0 & 0 & 0 & 0 & 0 & 0 & 0 & 0 & 0 & 0 & 0 & 0 & 0 & 0 & 0 & 0 & 1 & \alpha  & 1 & 0 \\
 0 & 0 & 0 & 0 & 0 & 0 & 0 & 0 & 0 & 0 & 0 & 0 & 0 & 0 & 0 & 0 & 0 & 0 & 0 & 1 & \alpha  & 1 \\
 0 & 0 & 0 & 0 & 0 & 0 & 0 & 0 & 0 & 0 & 0 & 0 & 0 & 0 & 0 & 0 & 0 & 0 & 0 & 0 & 1 & \alpha  \\
\end{array}
\right)\]

\bajm{[I want to write some determinants here too]}
\ignore{{   %% this is n=9 without the guide lines: 
\[\left(
\begin{array}{cccccccccccccccccccccc}
 \alpha  & 1 & 1 & 1 & 1 & 0 & 0 & 0 & 0 & 0 & 0 & 0 & 0 & 0 & 0 & 0 & 0 & 0 & 0 & 0 & 0 & 0 \\
 1 & \alpha  & 1 & 1 & 0 & 1 & 0 & 0 & 0 & 0 & 0 & 0 & 0 & 0 & 0 & 0 & 0 & 0 & 0 & 0 & 0 & 0 \\
 1 & 1 & \alpha  & 0 & 1 & 1 & -\frac{1}{2} & 0 & 0 & 0 & 0 & 0 & 0 & 0 & 0 & 0 & 0 & -\frac{\sqrt{3}}{2} & 0 & 0 & 0 & 0 \\
 1 & 1 & 0 & \alpha  & 1 & 1 & 0 & 0 & 0 & 0 & 0 & 0 & 0 & 0 & 0 & 0 & 0 & 0 & 0 & 0 & 0 & 0 \\
 1 & 0 & 1 & 1 & \alpha  & 1 & -\frac{1}{2} & 0 & 0 & 0 & 0 & 0 & 0 & 0 & 0 & 0 & 0 & \frac{\sqrt{3}}{2} & 0 & 0 & 0 & 0 \\
 0 & 1 & 1 & 1 & 1 & \alpha  & 1 & 0 & 0 & 0 & 0 & 0 & 0 & 0 & 0 & 0 & 0 & 0 & 0 & 0 & 0 & 0 \\
 0 & 0 & -\frac{1}{2} & 0 & -\frac{1}{2} & 1 & \alpha  & 1 & 0 & 0 & 0 & 0 & 0 & \frac{\sqrt{3}}{2} & 0 & \frac{\sqrt{3}}{2} & 0 & 0 & 0 & 0 & 0 & 0 \\
 0 & 0 & 0 & 0 & 0 & 0 & 1 & \alpha  & 1 & 0 & 0 & 0 & 0 & 0 & 0 & 0 & 0 & 0 & 0 & 0 & 0 & 0 \\
 0 & 0 & 0 & 0 & 0 & 0 & 0 & 1 & \alpha  & 1 & 0 & 0 & 0 & 0 & 0 & 0 & 0 & 0 & 0 & 0 & 0 & 0 \\
 0 & 0 & 0 & 0 & 0 & 0 & 0 & 0 & 1 & \alpha  & 1 & 0 & 0 & 0 & 0 & 0 & 0 & 0 & 0 & 0 & 0 & 0 \\
 0 & 0 & 0 & 0 & 0 & 0 & 0 & 0 & 0 & 1 & \alpha  & 0 & 0 & 0 & 0 & 0 & 0 & 0 & 0 & 0 & 0 & 0 \\
 0 & 0 & 0 & 0 & 0 & 0 & 0 & 0 & 0 & 0 & 0 & \alpha  & 1 & -1 & 1 & -1 & 0 & 0 & 0 & 0 & 0 & 0 \\
 0 & 0 & 0 & 0 & 0 & 0 & 0 & 0 & 0 & 0 & 0 & 1 & \alpha  & 1 & 1 & 0 & -1 & 0 & 0 & 0 & 0 & 0 \\
 0 & 0 & 0 & 0 & 0 & 0 & \frac{\sqrt{3}}{2} & 0 & 0 & 0 & 0 & -1 & 1 & \alpha  & 0 & 1 & -1 & -\frac{1}{2} & 0 & 0 & 0 & 0 \\
 0 & 0 & 0 & 0 & 0 & 0 & 0 & 0 & 0 & 0 & 0 & 1 & 1 & 0 & \alpha  & 1 & 1 & 0 & 0 & 0 & 0 & 0 \\
 0 & 0 & 0 & 0 & 0 & 0 & \frac{\sqrt{3}}{2} & 0 & 0 & 0 & 0 & -1 & 0 & 1 & 1 & \alpha  & 1 & \frac{1}{2} & 0 & 0 & 0 & 0 \\
 0 & 0 & 0 & 0 & 0 & 0 & 0 & 0 & 0 & 0 & 0 & 0 & -1 & -1 & 1 & 1 & \alpha  & 1 & 0 & 0 & 0 & 0 \\
 0 & 0 & -\frac{\sqrt{3}}{2} & 0 & \frac{\sqrt{3}}{2} & 0 & 0 & 0 & 0 & 0 & 0 & 0 & 0 & -\frac{1}{2} & 0 & \frac{1}{2} & 1 & \alpha  & 1 & 0 & 0 & 0 \\
 0 & 0 & 0 & 0 & 0 & 0 & 0 & 0 & 0 & 0 & 0 & 0 & 0 & 0 & 0 & 0 & 0 & 1 & \alpha  & 1 & 0 & 0 \\
 0 & 0 & 0 & 0 & 0 & 0 & 0 & 0 & 0 & 0 & 0 & 0 & 0 & 0 & 0 & 0 & 0 & 0 & 1 & \alpha  & 1 & 0 \\
 0 & 0 & 0 & 0 & 0 & 0 & 0 & 0 & 0 & 0 & 0 & 0 & 0 & 0 & 0 & 0 & 0 & 0 & 0 & 1 & \alpha  & 1 \\
 0 & 0 & 0 & 0 & 0 & 0 & 0 & 0 & 0 & 0 & 0 & 0 & 0 & 0 & 0 & 0 & 0 & 0 & 0 & 0 & 1 & \alpha  \\
\end{array}
\right)\]
}}}

